\section{Notation and Assumption Bundles}\label{app:min-notation}

\subsection{Manifesto Running Example}

The main text fixes one policy dimension, $d=\mathrm{econ}$. For a manifesto
span $x$, $f(x)$ denotes the economic-policy readout that a trusted coder or
calibrated rubric scorer would assign for the public-services versus taxation
rubric. The evidence sufficient for that readout is not directly observed.
Observed or learned objects are surrogates for it: prompted summaries, neural
states, quasi-sentence aggregates, or root scores. In the reported Manifesto
benchmark, the observed labels are root-level expert targets for sampled
manifestos; node-level observations enter only in the local-law audit or
quasi-sentence supervision setting. The discrepancy metric is $\metric$.

\subsection{Tree Objects}

$\Strings$ is the set of finite token sequences. Concatenation is
$u\concat v$. A document $x$ is partitioned into contiguous leaves
$(b_1,\ldots,b_k)$. A local summarizer $g:\Strings\rightsquigarrow\Strings$
stores $s_i=g(b_i)$ at leaves and
$s_u=g(s_{u_L}\concat s_{u_R})$ at internal nodes. The raw span under node $u$
is denoted $S(u)$ and is used for audits, not for deployed bottom-up
execution. The readout $f:\Strings\rightsquigarrow\Y$ scores the root summary.

Two strings are oracle-equivalent, $u\sim v$, when
$\metric(f(u),f(v))=0$. Context compatibility says $u\sim v$
continues to hold after inserting either string into the same left or right
context. The Manifesto reading is that substituting an economically equivalent
span should not change the economic score of the surrounding platform.

\subsection{Assumption Bundles}

The Preservation Stack is C1, C3, and context compatibility. C1 is leaf
sufficiency: $g(b)\sim b$. C3 is merge consistency: split-then-merge preserves
the raw union span's oracle value. The Recompression Stack adds C2:
$g(s)\sim s$ for stored summaries in the range of $g$. The Optimization Stack
adds oracle-measurability of downstream DPO/GRPO-style losses. The Certificate
Stack adds sampling, calibration, transport, and clipping assumptions for the
finite-sample audit.

\subsection{Granularity Vocabulary}

The paper uses paragraphs, sections, chapters, and books as reader-facing
granularity levels. Formally these are choices of leaf size and tree
partition. The guarantee is not tied to a linguistic unit; it is tied to
whether the chosen unit supports C1/C2/C3 for the target oracle. Root-level
corpus evaluation and node-level local-law auditing are therefore different
observation designs over the same task oracle.
